% ============================================================
% GX 339-4 HAWK-I Reduction — Overleaf writeup
% Paste this file into Overleaf and upload the figures listed
% in each \includegraphics command.
% ============================================================

\documentclass[12pt,a4paper]{article}

% --- Packages -----------------------------------------------
\usepackage[margin=2.5cm]{geometry}
\usepackage{amsmath}
\usepackage{graphicx}
\usepackage{booktabs}
\usepackage{hyperref}
\usepackage{natbib}
\usepackage{caption}
\usepackage{subcaption}
\usepackage{float}
\usepackage{lineno}
\usepackage{setspace}
\onehalfspacing

% --- Title --------------------------------------------------
\title{Near-Infrared Differential Photometry of GX~339-4:\\
       HAWK-I Ks-Band Reduction and ZOGY Difference Imaging}
\author{Eoin Trainor\\
        \small Research Masters, University College Cork\\
        \small Supervisor: Dr.\ Mark Kennedy}
\date{\today}

% ============================================================
\begin{document}
\maketitle
\tableofcontents
\newpage

% ============================================================
\section{Introduction}
\label{sec:intro}
% ============================================================

GX~339-4 is one of the most frequently studied black hole low-mass X-ray
binaries (LMXBs), having undergone numerous outburst--quiescence cycles since
its discovery \citep{Markert1973}. Despite decades of observation, the mass of
the central black hole remains poorly constrained, largely because the system
spends much of its time in an active or semi-active state that contaminates the
donor star's photometric signal with accretion disc and jet emission.

The most model-independent route to the black hole mass is the binary mass
function, derived from the radial velocity curve and the orbital inclination.
The inclination can be constrained by fitting the ellipsoidal modulation of the
donor star's near-infrared (NIR) lightcurve with a Roche-lobe geometry model
such as \textsc{icarus} \citep{Breton2012}. In the NIR Ks band ($\lambda
\approx 2.15\,\mu$m), the donor star's thermal emission dominates over the
accretion disc during quiescence, making this the preferred bandpass for such
studies.

Two field stars in close angular proximity to GX~339-4 complicate
straightforward aperture photometry of the target. We therefore adopt
\textbf{ZOGY proper image subtraction} \citep{Zackay2016} to isolate the
time-variable flux of GX~339-4 from its non-variable neighbours.

This document describes the full data reduction pipeline applied to 12
observing blocks (OBs) of HAWK-I Ks-band imaging acquired at the ESO Very
Large Telescope (VLT), covering the calibration, alignment, difference imaging,
and lightcurve extraction stages.

% ============================================================
\section{Observations}
\label{sec:obs}
% ============================================================

\subsection{Instrument and Setup}

Observations were obtained with HAWK-I (High Acuity Wide field K-band Imager)
mounted at the Nasmyth A focus of VLT UT4 (Yepun). HAWK-I employs four
Hawaii-2RG $2048 \times 2048$ pixel detectors arranged in a $2 \times 2$
mosaic with a total field of view of $7.5' \times 7.5'$ and a plate scale of
$0.106''$/pixel. GX~339-4 was imaged on the bottom-left chip (Detector 1) in
all observing blocks.

All science observations used the Ks broadband filter ($\lambda_c =
2.146\,\mu$m, $\Delta\lambda = 0.325\,\mu$m), with an integration time of
DIT~$= 10$\,s and NDIT~$= 9$ co-added sub-exposures per saved frame, giving
an effective on-sky time of 90\,s per FITS file.

\subsection{Observing Log}

Twelve observing blocks were executed between 2025 May 17 and 2025 September 20,
yielding 317 raw science frames. Table~\ref{tab:obs_log} summarises the
campaign.

\begin{table}[H]
\centering
\caption{Observing log. OBs~11--12 are excluded from all analysis
(confirmed outburst; see Section~\ref{sec:data_selection}).}
\label{tab:obs_log}
\begin{tabular}{lccrl}
\toprule
OB & Date (UT) & Frames & Exp.\ (s) & Notes \\
\midrule
GX339\_Ks\_1  & 2025-05-17 & 26 & 2340 & \\
GX339\_Ks\_2  & 2025-06-02 & 26 & 2340 & \\
GX339\_Ks\_3  & 2025-06-04 & 26 & 2340 & \\
GX339\_Ks\_4  & 2025-06-19 & 26 & 2340 & \\
GX339\_Ks\_5  & 2025-07-03 & 26 & 2340 & \\
GX339\_Ks\_6  & 2025-07-14 & 26 & 2340 & \\
GX339\_Ks\_7  & 2025-07-18 & 26 & 2340 & \\
GX339\_Ks\_8  & 2025-07-20 & 26 & 2340 & \\
GX339\_Ks\_9  & 2025-07-23 & 31 & 2790 & OB interrupted; split into two ESO ADPs \\
GX339\_Ks\_10 & 2025-08-15 & 26 & 2340 & Suspected non-quiescent \\
GX339\_Ks\_11 & 2025-09-20 & 26 & 2340 & Outburst (Swift/BAT confirmed) \\
GX339\_Ks\_12 & 2025-08-30 & 26 & 2340 & Outburst (Swift/BAT confirmed) \\
\midrule
Total & & 317 & $\sim$28530 & \\
\bottomrule
\end{tabular}
\end{table}

\subsection{Calibration Frames}

A total of 240 dark frames (EXPTIME~$= 10$\,s, NDIT~$= 9$) and 234 Ks-band
sky flat frames were obtained. Flat frames were split into two DIT groups
($1.68$\,s and $3.5$\,s); the DIT~$= 1.68$\,s group was selected for master
flat construction on the basis of higher signal-to-noise.

\subsection{Data Selection}
\label{sec:data_selection}

OBs~11 and 12 were obtained during a confirmed outburst episode, identified by
cross-matching the OB epochs against the Swift/BAT hard X-ray (15--50\,keV)
long-term light curve. These two OBs are excluded from all reduction and
analysis steps. OB~10 was provisionally retained as science frames but excluded
from the ZOGY reference coadd on suspicion of non-quiescent behaviour. The
final analysis therefore uses OBs~1--9 ($N = 236$ frames after alignment
quality cuts).

% ============================================================
\section{Data Reduction Pipeline}
\label{sec:reduction}
% ============================================================

\subsection{Master Dark and Bad Pixel Mask}

A master dark frame was constructed by median-combining all 240 dark frames
after sigma-clipped outlier rejection ($3\sigma$, 5 iterations). A bad pixel
mask was derived from pixels exceeding $8\sigma$ above the master dark RMS map,
and was applied in all subsequent calibration steps.

\subsection{Master Flat Field}

The 58 DIT~$= 1.68$\,s sky flat frames were median-combined and normalised to
unit median, producing a master Ks-band flat field. Pixels deviating by more
than 50\% from the median were flagged and added to the bad pixel mask.

\subsection{Science Frame Calibration}

All 317 science frames were processed as follows:
\begin{enumerate}
  \item Dark subtraction using the master dark (scaled to DIT if required).
  \item Flat-field division using the master flat.
  \item Sky background subtraction: the sky level was estimated from the
        sigma-clipped median of each frame and subtracted, yielding
        sky-subtracted, calibrated frames with a median sky residual $< 1$\,ADU.
\end{enumerate}

\subsection{Astrometric Alignment}

All 317 calibrated frames were aligned to a common pixel grid using
\texttt{astroalign} \citep{Beroiz2020}, which matches star triangles between
frames without requiring a WCS solution. A total of 314/317 frames were
successfully aligned (3 frames rejected due to insufficient stars or excessive
distortion). The aligned pixel scale is $0.106''$/px throughout.

% ============================================================
\section{ZOGY Proper Image Subtraction}
\label{sec:zogy}
% ============================================================

\subsection{Method}

We implemented the ZOGY algorithm of \citet{Zackay2016} for optimal image
subtraction in the Fourier domain. ZOGY produces a \emph{proper difference
image} $D$ that is statistically optimal in the sense that its noise properties
are spatially uniform and its pixels are uncorrelated, unlike classical
image-subtraction methods. The key equations are:

\begin{equation}
\hat{D} = \frac{F_r\,\hat{P}_r\,\hat{N} - F_n\,\hat{P}_n\,\hat{R}}
               {\sqrt{F_r^2\,\sigma_n^2\,|\hat{P}_r|^2 +
                      F_n^2\,\sigma_r^2\,|\hat{P}_n|^2}}
\label{eq:D}
\end{equation}

\begin{equation}
\hat{P}_D = \frac{F_r\,F_n\,\hat{P}_r\,\hat{P}_n}{F_D\,\hat{D}_{\rm denom}}
\label{eq:PD}
\end{equation}

\begin{equation}
F_D = \frac{F_r\,F_n}{\sqrt{F_r^2\,\sigma_n^2 + F_n^2\,\sigma_r^2}}
\label{eq:FD}
\end{equation}

where $\hat{X}$ denotes the discrete Fourier transform of $X$; $R$ and $N$ are
the reference and new (science) images respectively; $P_r$, $P_n$ are their
PSFs; $\sigma_r$, $\sigma_n$ are their background RMS values; and $F_r$, $F_n$
are their photometric flux normalisations.

The matched-filter detection statistic $S$ (Equation~17 of \citealt{Zackay2016}),
which gives the optimal signal-to-noise estimator for a point source, is:

\begin{equation}
\hat{S} = F_D\,\hat{P}_D^*\,\hat{D}
\label{eq:S}
\end{equation}

The differential flux of a variable point source at position $(x_0, y_0)$ is
recovered as $\Delta f = S(x_0,\,y_0)\,/\,F_D$.

\subsection{Reference Image Construction}

The coadded reference image $R$ was built from all 236 aligned frames from
OBs~1--9, using the inverse-variance weighted coaddition scheme of
\citet{ZackayOfek2015}:

\begin{equation}
R = \frac{\sum_j w_j\,I_j}{\sum_j w_j}, \qquad w_j = \frac{F_j}{\sigma_j^2}
\end{equation}

where $F_j$ and $\sigma_j$ are the photometric flux normalisation and
background RMS of the $j$-th frame, estimated per-frame via 2MASS Ks-band
catalogue matching with annular sky subtraction. The resulting reference image
is shown in Figure~\ref{fig:reference}.

\begin{figure}[H]
  \centering
  % Upload: fig1_reference_image.png from report_figures/
  \includegraphics[width=\textwidth]{fig1_reference_image.png}
  \caption{Coadded ZOGY reference image $R$ (236 frames, Ks band). Left:
           full $2048\times2048$ detector field with GX~339-4 marked (red
           circle). Right: $\pm80$\,px zoom on GX~339-4. Two contaminating
           field stars are visible in close proximity to the target.}
  \label{fig:reference}
\end{figure}

\subsection{Per-Frame PSF Modelling}

For each science frame the empirical PSF was estimated by:
\begin{enumerate}
  \item Source extraction with \textsc{sep} \citep{Barbary2016} at a
        $5\sigma$ detection threshold.
  \item Cross-matching detected sources with the 2MASS Ks catalogue
        \citep{Skrutskie2006} within a $7'$ search radius, retaining stars
        with $9.5 \leq K_s \leq 14.5$\,mag.
  \item Selecting up to 30 isolated, unsaturated stars (neighbour separation
        $> 6 \times$ FWHM) and median-combining their $51 \times 51$ pixel
        cutouts into a normalised PSF stamp.
  \item Fitting a \textbf{Moffat profile} to the normalised PSF stamp.
        The Moffat parameterisation was adopted over a Gaussian because it
        better reproduces the extended wings of seeing-broadened stellar
        PSFs, reducing residual power in the difference image $D$ at large
        separations from subtracted sources.  The fitted Moffat profile
        serves as $P_n$ in Equation~\ref{eq:D}.
\end{enumerate}

The per-frame seeing FWHM ranged from $3.5$--$8.9$\,px ($0.37''$--$0.94''$)
across the campaign, with a median of $6.00$\,px ($0.64''$).

\subsection{Photometric Normalisation}

Before applying Equation~\ref{eq:D}, each science frame was scaled to the
reference photometric system by multiplying by $F_r / F_n$. This ensures
$F_r = F_n^{\rm eff}$ in the subtraction, so that the source term
$S \cdot (F_r - F_n)$ vanishes and static sources cancel exactly.

\subsection{Difference Image Quality}

\begin{figure}[H]
  \centering
  % Upload: fig2_before_after.png from report_figures/
  \includegraphics[width=\textwidth]{fig2_before_after.png}
  \caption{ZOGY before/after comparison, $\pm80$\,px zoom around GX~339-4.
           Left: individual science frame $N$. Centre: difference image $D$,
           displayed on a symmetric colour scale ($\pm 5\sigma$). Right:
           reference $R$ (236-frame coadd). The two contaminating neighbours
           visible in $N$ and $R$ are suppressed in $D$.}
  \label{fig:before_after}
\end{figure}

The quality of the ZOGY subtraction was assessed using the sigma-clipped
standard deviation of $D$ in source-free background regions. In an ideal
subtraction, this quantity equals unity; values in the range $0.7$--$1.3$
are considered acceptable for a crowded Galactic field. For our 236 difference
images the median clipped $D_{\rm std} = 0.738$, an improvement over the
value of 0.726 obtained with Gaussian PSF modelling, and confirming
the noise model is well-calibrated (Figure~\ref{fig:dstd}).
The remaining deficit from the ideal value of 1.0 reflects spatially
varying PSF across the mosaic, sub-pixel centring errors, and
flat-fielding-induced correlated noise.

\begin{figure}[H]
  \centering
  % Upload: fig4_dstd_timeline.png from report_figures/
  \includegraphics[width=\textwidth]{fig4_dstd_timeline.png}
  \caption{Top: sigma-clipped $D$ standard deviation per frame across all
           nine OBs (Moffat PSF model).  Dashed lines mark the guideline
           range $0.7$--$1.3$; the majority of frames fall within this
           range, with a minority below 0.7 driven by frames with
           well-matched seeing.  Bottom: per-frame seeing FWHM from
           Moffat fits.  OBs are colour-coded.}
  \label{fig:dstd}
\end{figure}

\begin{figure}[H]
  \centering
  % Upload: fig5_d_histogram.png from report_figures/
  \includegraphics[width=0.65\textwidth]{fig5_d_histogram.png}
  \caption{Pixel value distribution of a representative difference image $D$
           (OB1, frame 1). The red curve shows a Gaussian fit centred on
           the sigma-clipped mean with $\sigma = 0.915$; the black dashed
           curve shows the ideal unit Gaussian ($\sigma = 1.0$) at the same
           centre. The measured width is consistent with the ZOGY noise model.
           The non-zero mean reflects a residual sky offset between $N$ and
           $R$, corrected locally at the target position in the photometry
           stage.}
  \label{fig:dhist}
\end{figure}

The $5\sigma$ point-source detection depth, estimated from the per-frame
background noise and PSF FWHM, is $K_s \approx 17.1$\,mag (8\,px aperture).

\subsection{GX~339-4 Position and Target Extraction}

The pixel position of GX~339-4 in the reference frame was determined by
seeding a \textsc{sep} centroid fit at the 2MASS catalogue position
(RA~$= 255.7058\deg$, Dec~$= -48.7898\deg$; J2000) with a constrained $8$\,px
search radius to avoid jumping to contaminating neighbours. The centroid
converged at pixel $(750,\,1012)$ (0-indexed).

The ZOGY matched-filter statistic $S$ (Equation~\ref{eq:S}) was evaluated at
this pixel position for each of the 236 frames. Eight photometric reference
stars from the 2MASS catalogue ($9.5 \leq K_s \leq 11.7$\,mag) were
extracted from $S$ at the same time as a sanity check; their per-frame $S$
values are consistent with zero (mean $< 1\%$ of the GX~339-4 scatter),
confirming that the subtraction is working correctly for non-variable sources.

% ============================================================
\section{Lightcurve Analysis}
\label{sec:lightcurve}
% ============================================================

\subsection{Orbital Ephemeris}

The orbital phase $\phi$ for each frame was computed using the ephemeris of
\citet{Heida2017}:

\begin{equation}
\phi = \frac{{\rm MJD} - T_0}{P} \pmod{1}, \qquad
T_0 = 57529.397\,{\rm MJD}, \quad P = 1.7587\,{\rm d}
\label{eq:phase}
\end{equation}

where $T_0$ is the time of inferior conjunction of the donor star (donor
closest to the observer). For ellipsoidal modulation, minima are expected at
$\phi = 0.0$ and $0.5$ (donor face-on, smallest projected area), and maxima at
$\phi = 0.25$ and $0.75$ (donor side-on, largest projected area).

\subsection{Raw Differential Lightcurve}

Figure~\ref{fig:lc_raw} shows the per-OB mean ZOGY $S_{\rm target}$ as a
function of time (MJD). A clear, monotonic brightening trend is apparent over
the 66-day campaign, with the differential flux rising by a factor of
approximately 4 between OB~1 (2025 May 17) and OB~7 (2025 Jul 18). This trend
is confirmed as real by the reference stars, which remain flat throughout.

\begin{figure}[H]
  \centering
  % Upload: lc06_ob_means_vs_phase.png from lightcurve/
  \includegraphics[width=\textwidth]{lc06_ob_means_vs_phase.png}
  \caption{Left: per-OB mean $S_{\rm target}$ as a function of MJD.
           Error bars are the standard error on the mean within each OB.
           A monotonic brightening of GX~339-4 is evident over the 66-day
           campaign. Right: the same 9 data points phase-folded on the
           Heida+2017 ephemeris ($P = 1.7587$\,d). Blue and red dotted lines
           mark the expected ellipsoidal minima and maxima respectively.
           Phase coverage is sparse, with gaps at $\phi = 0$--$0.29$ and
           $0.4$--$0.55$.}
  \label{fig:lc_raw}
\end{figure}

\subsection{Phase-Folded Lightcurve}

The 9 OBs sample 9 distinct orbital phases in the range $\phi \approx
0.29$--$0.94$. The right panel of Figure~\ref{fig:lc_raw} shows the
phase-folded per-OB means. No clear ellipsoidal modulation is discernible:
the inter-OB scatter is dominated by the long-term flux trend rather than the
expected $\sim 5$--$15\%$ peak-to-peak ellipsoidal amplitude. In particular,
OBs at similar orbital phases (e.g.\ OBs~1 and 9 at $\phi \approx 0.62$;
OBs~4 and 5 at $\phi \approx 0.30$) show significantly different $S$ values,
consistent with temporal variability rather than orbital modulation.

OB~7 (2025 Jul 18, $\phi = 0.83$) is anomalously bright relative to adjacent
OBs (factor $\sim 10$ above OB~6 at $\phi = 0.59$), and may indicate the
onset of enhanced accretion activity preceding the confirmed outburst in
OBs~11--12.

\subsection{Interpretation}

The monotonic brightening observed over OBs~1--9 indicates that GX~339-4 was
already entering a more active accretion state during the observing campaign.
The system was not in a stable quiescent state during this period, as
implicitly assumed by the ellipsoidal modulation analysis. Consequently:

\begin{itemize}
  \item The long-term flux trend ($\sim$ factor 4 over 66 days) is larger than
        the expected ellipsoidal signal ($\sim 5$--$15\%$) by more than an
        order of magnitude.
  \item The two signals cannot be separated with only 9 orbital phase
        measurements and no independent constraint on the trend shape.
  \item Extraction of the binary mass function from this dataset is not
        currently feasible.
\end{itemize}

% ============================================================
\section{Discussion and Future Work}
\label{sec:discussion}
% ============================================================

The ZOGY pipeline performs correctly: the difference image noise is well
characterised ($D_{\rm std}^{\rm clip} = 0.726$), reference stars are flat,
and the $5\sigma$ depth of $K_s \approx 17.1$\,mag is adequate for detecting
GX~339-4's quiescent flux ($K_s \approx 16.5$--$17$\,mag). The limiting factor
is the astrophysical state of the target during the campaign.

To successfully extract the ellipsoidal lightcurve and constrain the black hole
mass, future observations should:

\begin{enumerate}
  \item \textbf{Confirm quiescence} using Swift/BAT or MAXI X-ray monitoring
        before triggering NIR observations, and require a stable, low X-ray
        flux for several weeks prior.
  \item \textbf{Observe over 3--5 consecutive nights}, giving $\sim 1.7$--$2.8$
        full orbital cycles in a single, short-baseline campaign. This maximises
        phase coverage while minimising the risk of long-term variability
        building up.
  \item \textbf{Schedule OBs at target orbital phases} using the
        \citet{Heida2017} ephemeris to ensure uniform coverage across $\phi =
        0$--$1$, particularly at the currently unsampled phases
        $\phi = 0$--$0.29$ and $0.4$--$0.55$.
\end{enumerate}

% ============================================================
\section{Conclusions}
\label{sec:conclusions}
% ============================================================

We have presented a complete Ks-band imaging reduction pipeline for 12 HAWK-I
observing blocks of GX~339-4, incorporating dark subtraction, flat-fielding,
sky subtraction, astrometric alignment, and ZOGY proper image subtraction. The
pipeline successfully constructs a 236-frame coadded reference image, produces
per-frame difference images with well-characterised noise ($D_{\rm std}^{\rm
clip} = 0.726$), and extracts the ZOGY matched-filter flux of GX~339-4 and 8
photometric reference stars at each epoch.

Phase-folding the differential flux on the orbital ephemeris of
\citet{Heida2017} reveals that GX~339-4 was undergoing a monotonic brightening
over the 66-day campaign, indicative of rising accretion activity ahead of the
confirmed outburst in OBs~11--12. The ellipsoidal modulation signal is
swamped by this long-term trend and cannot be extracted from the current
dataset. New observations obtained during a confirmed quiescent epoch, covering
at least two full orbital periods, are required to proceed with ellipsoidal
lightcurve modelling and black hole mass determination.

% ============================================================
% References
% ============================================================
\begin{thebibliography}{99}

\bibitem[Barbary(2016)]{Barbary2016}
Barbary K., 2016, \textit{sep}: Source Extraction and Photometry.
JOSS, 1, 58

\bibitem[Beroiz et al.(2020)]{Beroiz2020}
Beroiz M., Cabral J.~B., Sanchez B., 2020, Astron.\ Comput., 32, 100384

\bibitem[Breton et al.(2012)]{Breton2012}
Breton R.~P. et al., 2012, ApJ, 769, 108

\bibitem[Heida et al.(2017)]{Heida2017}
Heida M. et al., 2017, ApJ, 846, 132.
$T_0 = 57529.397$\,MJD, $P = 1.7587$\,d (inferior conjunction of donor).

\bibitem[Markert et al.(1973)]{Markert1973}
Markert T.~H. et al., 1973, ApJ, 184, L67

\bibitem[Skrutskie et al.(2006)]{Skrutskie2006}
Skrutskie M.~F. et al., 2006, AJ, 131, 1163

\bibitem[Zackay \& Ofek(2015)]{ZackayOfek2015}
Zackay B., Ofek E.~O., 2017, ApJ, 836, 187

\bibitem[Zackay et al.(2016)]{Zackay2016}
Zackay B., Ofek E.~O., Gal-Yam A., 2016, ApJ, 830, 27

\end{thebibliography}

\end{document}
